\documentclass[11pt]{article}

\usepackage{amsmath}
\usepackage{amssymb}
\usepackage{booktabs}
\usepackage[margin=1in]{geometry}
\usepackage{hyperref}

\newcommand{\sg}{\operatorname{sg}}
\DeclareMathOperator{\Unif}{\mathcal{U}}

\title{Dreamer 4: Technical Torso\\ \large Decisions, Architectures, and Equations}
\author{(Extracted from Hafner, Yan, Lillicrap, ``Training Agents Inside of Scalable World Models'', 2025)}
\date{}

\begin{document}
\maketitle

This is a stripped-down skeleton: only the parts that actually carry information
(method decisions, architecture, equations, key configs). Motivation, related work,
and references are omitted.

%==================================================================
\section{Background}
%==================================================================

\subsection{Flow matching}
The network $f_\theta$ is trained to restore a clean datapoint $x_1$ from a corrupted
version $x_\tau$, where the signal level $\tau \in [0,1]$ mixes noise and data
($\tau=0$ pure noise, $\tau=1$ clean). Flow matching has the network predict the
velocity $v = x_1 - x_0$:
\begin{align}
  x_\tau &= (1-\tau)\,x_0 + \tau\,x_1,
  &x_0 &\sim \mathcal{N}(0,\mathbb{I}),
  &x_1 &\sim \mathcal{D},
  &\tau &\sim p(\tau), \\
  \mathcal{L}(\theta) &= \big\| f_\theta(x_\tau,\tau) - (x_1 - x_0) \big\|^2 .
\end{align}
$\tau$ is typically drawn uniform or logit-normal. Inference starts from noise $x_0$
and integrates over $K$ steps of size $d = 1/K$:
\begin{equation}
  x_{\tau+d} = x_\tau + f_\theta(x_\tau,\tau)\,d, \qquad x_0 \sim \mathcal{N}(0,\mathbb{I}).
\end{equation}

\subsection{Shortcut models}
Condition the network on the step size $d$ as well as $\tau$, so step size can be
chosen at inference. For the finest step $d_{\min}$ the flow loss is used; for larger
steps a bootstrap loss distills two half-steps:
\begin{align}
  x_0 &\sim \mathcal{N}(0,I), & x_1 &\sim \mathcal{D}, & \tau,d &\sim p(\tau,d), \\
  b' &= f_\theta(x_\tau,\tau,\tfrac{d}{2}), &
  b'' &= f_\theta(x',\tau+\tfrac{d}{2},\tfrac{d}{2}), &
  x' &= x_\tau + b'\tfrac{d}{2}, \\
  \mathcal{L}(\theta) &= \big\| f_\theta(x_\tau,\tau,d) - v_{\text{target}} \big\|^2,
  &v_{\text{target}} &=
  \begin{cases}
    x_1 - x_0 & \text{if } d = d_{\min}, \\
    \sg(b' + b'')/2 & \text{else.}
  \end{cases}
\end{align}
Step size sampled as a power of two; signal level uniform over the reachable grid:
\begin{equation}
  d \sim 1/\Unif(\{1,2,4,8,\dots,K_{\max}\}), \qquad
  \tau \sim \Unif(\{0,\, 1/d,\, \dots,\, 1 - 1/d\}).
\end{equation}
In practice $2$--$4$ sampling steps suffice versus $64+$ for ordinary diffusion.

\subsection{Diffusion forcing}
For sequences, assign a \emph{different} signal level to each time step. Every step
then doubles as a denoising target and as (noised) history context for later steps,
and inference supports flexible noise patterns (e.g.\ clean/lightly-noised history).

%==================================================================
\section{World Model Agent}
%==================================================================

The agent is a \textbf{causal tokenizer} (continuous video$\to$latent) plus an
\textbf{interactive dynamics model}, both built on the same block-causal transformer.
Tokenizer: masked-autoencoding reconstruction. Dynamics: shortcut-forcing objective,
giving fast interactive rollout with few forward passes and no error accumulation.

\paragraph{Three phases.}
\begin{enumerate}
  \item \textbf{World model pretraining} --- train tokenizer on videos (Eq.~9), then
    train dynamics on tokenized videos (and optionally actions) via shortcut forcing
    (Eq.~13).
  \item \textbf{Agent finetuning} --- insert task tokens; train policy + reward heads
    (Eqs.~13, 15).
  \item \textbf{Imagination training} --- optimize policy (Eq.~17) and value (Eq.~16)
    on world-model rollouts; transformer kept frozen.
\end{enumerate}
A single dynamics transformer carries multiple modalities and heads; all loss terms
are normalized by running RMS estimates.

%==================================================================
\section{Causal Tokenizer}
%==================================================================

Encoder--decoder with a low-dimensional bottleneck, causal in time (temporal
compression while still decoding frame-by-frame for interactive inference).

\paragraph{Architecture decisions.}
\begin{itemize}
  \item Each step = image patch tokens + learned latent tokens.
  \item After the encoder, read latents out via linear projection to a smaller channel
    dim followed by $\tanh$. Decoder projects back up, concatenated with learned tokens
    to read out patches.
  \item Cross-modal routing: in the encoder, latents attend to all modalities but each
    modality attends only within itself; in the decoder each modality attends within
    itself and to latents, latents only within themselves.
\end{itemize}

\paragraph{Masked autoencoding loss.}
\begin{equation}
  \mathcal{L}(\theta) = \mathcal{L}_{\text{MSE}}(\theta) + 0.2\,\mathcal{L}_{\text{LPIPS}}(\theta).
\end{equation}
Input patches are dropped with probability $p \sim U(0,0.9)$ (sampled per image,
including the $p=0$ inference case). MAE training improves spatial consistency of the
dynamics model's generations.

%==================================================================
\section{Interactive Dynamics}
%==================================================================

Operates on the interleaved sequence of actions and (frozen) tokenizer representations.
Trained with \textbf{shortcut forcing} for fast inference at $K=4$ forward passes per
generated frame.

\paragraph{Input tokenization.}
\begin{itemize}
  \item Representations linearly projected into $S_z$ spatial tokens, concatenated with
    $S_r$ learned register tokens and a single token encoding signal level + step size
    (both discrete $\to$ embedding lookup, channels concatenated).
  \item Each action component encoded into $S_a$ tokens and summed with a learned
    embedding (continuous: linear; categorical/binary: embedding lookup). For unlabeled
    video, only the learned embedding is used.
\end{itemize}

\paragraph{Shortcut forcing objective.}
Inputs are actions $a=\{a_t\}$, per-step signal levels $\tau=\{\tau_t\}$, step sizes
$d=\{d_t\}$, and corrupted representations $\tilde z = \{z_t^{(\tau_t)}\}$; the model
predicts clean representations $z_1 = \{z_t^1\}$. Here $t\in[1,T]$ is the sequence step
and $\tau_t\in[0,1]$ its signal level.
\begin{align}
  z_0 &\sim \mathcal{N}(0,1), & z_1 &\sim \mathcal{D}, & \tau,d &\sim p(\tau,d),
  & \tau,d &\in [0,1]^T, \\
  \hat z_1 &= f_\theta(\tilde z,\tau,d,a), & \tilde z &= (1-\tau)\,z_0 + \tau\,z_1.
\end{align}

Key decision: predict \textbf{clean representations (x-prediction)}, not velocities
(v-prediction). V-prediction trains high-frequency outputs whose subtle errors
accumulate over long frame-by-frame rollouts; x-prediction gives stable rollouts of
arbitrary length. The flow term is computed directly in x-space; the bootstrap term is
formed in v-space and scaled back to x-space:
\begin{align}
  b' &= \big(f_\theta(\tilde z,\tau,\tfrac{d}{2},a) - z_\tau\big)\big/(1-\tau),
  & z' &= \tilde z + b'\tfrac{d}{2}, \\
  b'' &= \big(f_\theta(z',\tau+\tfrac{d}{2},\tfrac{d}{2},a) - z'\big)\big/\big(1-(\tau+\tfrac{d}{2})\big), \\
  \mathcal{L}(\theta) &=
  \begin{cases}
    \big\|\hat z_1 - z_1\big\|_2^2 & \text{if } d = d_{\min}, \\[4pt]
    (1-\tau)^2\,\Big\| (\hat z_1 - \tilde z)/(1-\tau) - \sg(b' + b'')/2 \Big\|_2^2 & \text{else.}
  \end{cases}
\end{align}
(The $(1-\tau)^2$ factor brings the v-space bootstrap loss into the same range as the
x-space flow loss, via $\|\hat x_1 - x_1\|_2^2 = (1-\tau)^2\|\hat v_\tau - v_\tau\|_2^2$
with $\hat v_\tau = (\hat x_1 - x_\tau)/(1-\tau)$.)

\paragraph{Ramp loss weight.}
Low signal levels carry little learning signal (the flow term degenerates to the
dataset mean) while the bootstrap term is easier (deterministic targets). Up-weight
high signal levels linearly:
\begin{equation}
  w(\tau) = 0.9\,\tau + 0.1.
\end{equation}

\paragraph{Inference.}
Autoregressive in time; each frame generated with the shortcut model at $K=4$ steps,
$d=1/4$. Past context inputs are slightly corrupted to $\tau_{\text{ctx}} = 0.1$ for
robustness to imperfect self-generations.

%==================================================================
\section{Imagination Training}
%==================================================================

\paragraph{Agent tokens.}
Task embeddings enter as an extra modality interleaved with image reps, actions, and
register tokens. Agent tokens attend to themselves and all other modalities, but
\emph{no modality attends back to them} --- crucial to avoid causal confusion: the
world model's future may only be influenced by actions, not by the current task.

\paragraph{Behavior cloning + reward model (multi-token prediction, $L=8$).}
\begin{equation}
  \mathcal{L}(\theta) = -\sum_{n=0}^{L}\ln p_\theta(a_{t+n}\mid h_t)
                        -\sum_{n=0}^{L}\ln p_\theta(r_{t+n}\mid h_t).
\end{equation}
Pretraining setting is reused (noisy reps, video-prediction loss kept) to preserve
capabilities. Policy/reward heads = small MLPs, one output layer per MTP distance.
Reward head: symexp twohot. Policy head: categorical or vectorized-binary.

\paragraph{Reinforcement learning in imagination.}
Initialize a value head and a frozen copy of the policy as a behavioral prior; update
only policy + value heads, transformer frozen. One rollout per dataset context
(diversity over depth). Rollouts unroll the transformer with itself, sampling reps from
the flow head and actions from the policy head, then annotate rewards and values.

\paragraph{Value head (TD-learning, $\lambda$-returns, $\gamma=0.997$).}
\begin{equation}
  \mathcal{L}(\theta) = -\sum_{t=1}^{T}\ln p_\theta(R_t^\lambda \mid s_t),
  \quad
  R_t^\lambda = r_t + \gamma c_t\big((1-\lambda)v_t + \lambda R_{t+1}^\lambda\big),
  \quad
  R_T^\lambda = v_T,
\end{equation}
with $c_t$ a non-terminal indicator; symexp twohot output.

\paragraph{Policy (PMPO).}
Uses only the \emph{sign} of advantages $A_t = R_t^\lambda - v_t$ (no return/advantage
normalization needed; equal focus across tasks). Split states into
$\mathcal{D}_+ = \{s_i \mid A_t \ge 0\}$ and $\mathcal{D}_- = \{s_i \mid A_t < 0\}$:
\begin{equation}
  \mathcal{L}(\theta) =
  \frac{1-\alpha}{|\mathcal{D}_-|}\sum_{i\in\mathcal{D}_-}\ln\pi_\theta(a_i\mid s_i)
  - \frac{\alpha}{|\mathcal{D}_+|}\sum_{i\in\mathcal{D}_+}\ln\pi_\theta(a_i\mid s_i)
  + \frac{\beta}{N}\sum_{i=1}^{N} \mathrm{KL}\big[\pi_\theta(a_i\mid s_i)\,\big\|\,\pi_{\text{prior}}\big].
\end{equation}
$\alpha = 0.5$ (balance positive/negative sets), $\beta = 0.3$ (behavioral prior).
The prior KL uses a \emph{reverse} direction vs.\ original PMPO, to keep the policy
inside the space of reasonable behaviors. All three terms are in nats, making the
scaling robust.

%==================================================================
\section{Efficient Transformer}
%==================================================================

Shared by tokenizer and dynamics. A 2D (time$\times$space) transformer, attention masked
\emph{causal in time} (all tokens in a step attend to each other and to the past).
Base: pre-layer RMSNorm, RoPE, SwiGLU; plus QKNorm and attention-logit soft capping for
stability.

\paragraph{Efficiency decisions.}
\begin{itemize}
  \item Separate space-only and time-only attention layers (break up dense
    all-token attention).
  \item Temporal attention only \emph{once every 4 layers} --- speeds up train+inference
    and improves quality (spatial inductive bias).
  \item GQA on all dynamics attention layers (shrink KV cache).
\end{itemize}

\paragraph{Sequence length.}
More spatial tokens $\to$ visual quality; more temporal tokens $\to$ longer consistent
context. Train on alternating short/long batches, then finetune on long batches only.
Batch length must exceed context length so the model doesn't overfit to always seeing a
start frame, enabling length generalization.

%==================================================================
\section{Key Configs and Headline Results}
%==================================================================

\paragraph{Model / compute.}
2B params total (400M tokenizer, 1.6B dynamics). Trained on 256--1024 TPU-v5p, batch
size 1/device, FSDP. 30\% of batch videos treated as standalone images (to learn start
frames). Minecraft: 256 spatial tokens, 192-frame context, 256 batch length. Real-world
datasets: 512 spatial tokens, 96-frame context, 128 batch length.

\paragraph{Data.}
VPT contractor dataset: 2541 h, 360p, mouse+keyboard at 20\,FPS. Keyboard = 23 binary
dists; mouse = categorical with 121 classes (foveated discretization, $11\times11$).
Image 360$\times$640, zero-padded to 384$\times$640, patch $16\times16 \to 960$ tokens;
tokenizer bottleneck $512\times16$ reshaped to $256\times32$.

\paragraph{Offline diamond challenge (1000 episodes, 60 min, empty inventory).}
Dreamer 4 is the first to obtain diamonds purely offline: $0.7\%$ diamond success,
$29\%$ iron pickaxe, $>90\%$ up to stone pickaxe. Beats OpenAI VPT (finetuned) using
$100\times$ less data, and nearly triples the VLA (Gemma 3) iron-pickaxe rate.
Imagination training helps more the harder the milestone, and reaches milestones faster.

\paragraph{World-model comparison (single H100).}
2B params, 640$\times$360, \textbf{9.6\,s context} (6$\times$ prior), 21\,FPS (real-time),
$14/16$ interaction tasks --- vs.\ Oasis-large $5/16$ ($\sim$5\,FPS) and Lucid-v1 $0/16$.

\paragraph{Action generalization.}
With only 10\,h of paired actions (out of 2541\,h video): 53\% PSNR / 75\% SSIM relative
to all-actions; with 100\,h: 85\% PSNR / 100\% SSIM. Action conditioning learned only on
the Overworld extrapolates to Nether/End (seen only as unlabeled video): 76\% PSNR,
80\% SSIM.

\paragraph{Ablation cascade (FVD$\downarrow$, train s/step, inference FPS).}
\begin{center}
\begin{tabular}{lccc}
\toprule
Model & Train s & FPS & FVD \\
\midrule
Diffusion Forcing Transformer        & 9.8 & 0.8  & 306 \\
\quad + Fewer sampling steps ($K{=}4$) & 9.8 & 9.1  & 875 \\
\quad + Shortcut model                & 9.8 & 9.1  & 329 \\
\quad + X-Prediction                  & 9.8 & 9.1  & 326 \\
\quad + X-Loss                        & 9.8 & 9.1  & 151 \\
\quad + Ramp weight                   & 9.8 & 9.1  & 102 \\
\quad + Alternating batch lengths     & 1.5 & 9.1  & 80  \\
\quad + Long context every 4 layers   & 0.6 & 18.9 & 70  \\
\quad + GQA                           & 0.5 & 23.2 & 71  \\
\quad + Time-factorized long context  & 0.4 & 30.1 & 91  \\
\quad + Register tokens               & 0.5 & 28.9 & 91  \\
\quad + More spatial tokens ($N_z{=}128$) & 0.8 & 25.7 & 66  \\
\quad + More spatial tokens ($N_z{=}256$) & 1.7 & 21.4 & 57  \\
\bottomrule
\end{tabular}
\end{center}
Final FVD 57 vs.\ 306 naive baseline (and 124 for the full architecture with v-space
prediction/losses). Shortcut forcing at 4 steps $\approx$ diffusion forcing at 64 steps
($16\times$ faster).

\end{document}
